\scanpage{171}

U N I V E R S I T E

DES SCIENCES ET TECHNIQUES

DU LANGUEDOC

\verb|^^^^^^^^^^^^|

\hfill Montpellier, le 14 octobre 1980

Service de la Scolarité

MA/gs

Diffusion : Liste A

N° 5254

\begin{center}
\uline{AVIS DE SOUTENANCE D'UNE THESE DE DOCTEUR-INGENIEUR}

``\uline{AGRONOMIE OPTION ZOOTECHNIE}''
\end{center}

\uline{NOM et prénom du candidat} : DELAVEAU Alain

\uline{Titre de la thèse} : ``CONTRIBUTION A L'ETUDE DE LA MORTALITE DES LAPEREAUX ENTRE LA NAISSANCE ET LE SEVRAGE - (Aspects\unsure{.} physio-\newline
logiques et zootechniques).''

\uline{Date, heure et lieu de soutenance prévus} : 17 octobre 1980 à 9 heures,30

Ecole Nationale Supérieure Agronomique

\uline{Nom du directeur de thèse} : M. L. DAUZIER

\uline{Désignation du laboratoire où la thèse est préparée} : Laboratoire de Physiologie animale à l'E.N.S.A. de Montpellier

\uline{Analyse de l'ouvrage} : La mortalité des lapereaux entre la naissance et le sevrage n'apparaît pas comme un phénomène aléatoire puisque l'essentiel des pertes seproduit durant la première semaine après la mise bas. On note d'ailleurs des causes identiques de pertes de lapereaux dans des élevages expérimentaux différents.

La qualité du nid joue, de plus, un rôle essentiel. C'est ainsi qu'il est possible de diminuer la mortalité post-natale en aménageant le nid le jour, ou même, 10 jours après la parturition. De plus, nous avons enregistré moins de pertes en utilisant des boîtes à nid équipées d'un fond creux. Mais, nous n'avons pas noté d'effet de la température interne de cette même boîte à nid sur les pertes de lapereaux.

Enfin, nous n'avons pas mis en évidence d'incidence de la séparation mère-lapereaux sur la mortalité entre la naissance et le sevrage.

\uline{Durée approximative des recherches ayant abouti à la thèse} : 4 ans

\uline{Composition du jury proposé} :
\begin{tabular}{l}
M. L. DAUZIER\\
M. JD. BAYLE\\
M. M. PRUD'HON\\
M. F. LEBAS\\
M. Y. REYNES
\end{tabular}

\begin{flushright}
P/ Le Secrétaire Général\\
L'Attachée d'administration\\
Universitaire,

\unsure{Gentil}

\uline{A. GENTIL}
\end{flushright}
\margin{171}
